% Created: 2026-09-30 04:37 -- Memo mathematique des isometries de Pentapol.
% docs/MEMO_ISOMETRIES_TERMINALE.tex
\documentclass[11pt,a4paper]{article}

\usepackage[T1]{fontenc}
\usepackage[utf8]{inputenc}
\usepackage[french]{babel}
\usepackage{amsmath,amssymb}
\usepackage{booktabs}
\usepackage{array}
\usepackage{geometry}
\usepackage{enumitem}
\usepackage{xcolor}
\usepackage{graphicx}
\usepackage{hyperref}
\usepackage{fancyhdr}

\geometry{margin=2cm}
\hypersetup{colorlinks=true,linkcolor=blue!55!black,urlcolor=blue!55!black}
\setlist[itemize]{leftmargin=1.5em,itemsep=0.25em}
\setlist[enumerate]{leftmargin=1.7em,itemsep=0.35em}
\renewcommand{\arraystretch}{1.25}
\pagestyle{fancy}
\fancyhf{}
\lhead{Pentapol}
\rhead{Isom\'etries et exponentielle complexe}
\cfoot{\thepage}

\newcommand{\C}{\mathbb{C}}
\newcommand{\Z}{\mathbb{Z}}
\newcommand{\id}{\mathrm{Id}}
\newcommand{\conj}[1]{\overline{#1}}

\title{\textbf{Visualiser les isom\'etries de Pentapol}\\
\large M\'emo de math\'ematiques pour un \'el\`eve de Terminale}
\author{}
\date{30 septembre 2026}

\begin{document}
\maketitle

\begin{center}
\fbox{\parbox{0.9\linewidth}{
\textbf{Id\'ee directrice.} Les rotations de Pentapol sont les puissances d'un quart de tour.
Dans le plan complexe, elles s'\'ecrivent avec $e^{ik\pi/2}$. Les sym\'etries axiales demandent
en plus la conjugaison complexe $z\mapsto\conj z$.
}}
\end{center}

\section{Les quatre gestes \'el\'ementaires}

Pentapol propose quatre transformations directes d'une pi\`ece :

\begin{itemize}
  \item une rotation d'un quart de tour dans le sens trigonom\'etrique ;
  \item une rotation d'un quart de tour dans le sens horaire ;
  \item une sym\'etrie axiale horizontale ;
  \item une sym\'etrie axiale verticale.
\end{itemize}

Nous utilisons le rep\`ere math\'ematique habituel : l'axe des abscisses est dirig\'e vers la
droite et l'axe des ordonn\'ees vers le haut. L'\'ecran d'un t\'el\'ephone utilise souvent un axe
vertical dirig\'e vers le bas ; les relations math\'ematiques restent cependant identiques.

\section{Repr\'esenter une pi\`ece par des nombres complexes}

Un point $M(x,y)$ est repr\'esent\'e par le nombre complexe
\[
  z=x+iy.
\]
On place d'abord le centre $O$ de la transformation \`a l'origine. Si le centre r\'eel a pour
affixe $c$, la formule compl\`ete s'obtient par
\[
  z'=c+T(z-c),
\]
o\`u $T$ d\'esigne la transformation effectu\'ee autour de l'origine.

\section{L'exponentielle complexe et les rotations}

La formule d'Euler est
\[
  e^{i\theta}=\cos\theta+i\sin\theta.
\]
Multiplier $z$ par $e^{i\theta}$ fait tourner le point d'un angle $\theta$ autour de l'origine :
\[
  R_{\theta}(z)=e^{i\theta}z.
\]

Dans Pentapol, les angles sont uniquement des multiples de $\pi/2$. Posons
\[
  r=R_{\pi/2},
  \qquad r(z)=e^{i\pi/2}z=iz.
\]
La transformation $r$ est le quart de tour dans le sens trigonom\'etrique. Ses puissances donnent
toutes les rotations disponibles :
\[
  r^k(z)=e^{ik\pi/2}z=i^kz.
\]

\subsection*{Le cycle des puissances de $i$}

\[
\begin{array}{ccccc}
1 & \longrightarrow & i & \longrightarrow & -1\\
\rotatebox[origin=c]{90}{$\longleftarrow$} &&&& \rotatebox[origin=c]{90}{$\longrightarrow$}\\[-0.4em]
&& -i &&
\end{array}
\qquad\text{puis retour \`a }1.
\]

\begin{center}
\begin{tabular}{ccccc}
\toprule
Puissance & Exponentielle & Facteur & Angle & $(x,y)\mapsto$\\
\midrule
$r^0$ & $e^{0}$ & $1$ & $0$ & $(x,y)$\\
$r^1$ & $e^{i\pi/2}$ & $i$ & $\pi/2$ & $(-y,x)$\\
$r^2$ & $e^{i\pi}$ & $-1$ & $\pi$ & $(-x,-y)$\\
$r^3$ & $e^{i3\pi/2}$ & $-i$ & $3\pi/2$ & $(y,-x)$\\
$r^4$ & $e^{i2\pi}$ & $1$ & $2\pi$ & $(x,y)$\\
\bottomrule
\end{tabular}
\end{center}

Les exposants sont calcul\'es modulo $4$ :
\[
  r^5=r,qquad r^7=r^3,qquad r^{10}=r^2,qquad r^{-1}=r^3.
\]

\section{La conjugaison et les sym\'etries axiales}

Une multiplication par $e^{i\theta}$ conserve l'orientation : elle produit une rotation, jamais
une r\'eflexion. Pour obtenir une sym\'etrie, il faut utiliser la conjugaison complexe :
\[
  z=x+iy
  \qquad\Longrightarrow\qquad
  \conj z=x-iy.
\]

La transformation
\[
  s(z)=\conj z
\]
est la sym\'etrie d'axe horizontal. Elle v\'erifie
\[
  s^2=\id.
\]

Plus g\'en\'eralement, la sym\'etrie dont l'axe forme un angle $\varphi$ avec l'axe horizontal
s'\'ecrit
\[
  \boxed{\;S_{\varphi}(z)=e^{2i\varphi}\conj z\;}.
\]

\section{Les quatre sym\'etries du carr\'e}

Toutes les r\'eflexions s'\'ecrivent sous la forme $r^ks$, avec $k\in\Z$. Par convention, dans
un produit comme $rs$, l'op\'eration de droite est effectu\'ee la premi\`ere : on applique $s$,
puis $r$.

\begin{center}
\begin{tabular}{cccc}
\toprule
Notation & Formule complexe & $(x,y)\mapsto$ & Axe\\
\midrule
$s$ & $\conj z$ & $(x,-y)$ & horizontal\\
$rs$ & $i\conj z$ & $(y,x)$ & $y=x$\\
$r^2s$ & $-\conj z$ & $(-x,y)$ & vertical\\
$r^3s$ & $-i\conj z$ & $(-y,-x)$ & $y=-x$\\
\bottomrule
\end{tabular}
\end{center}

Pentapol poss\`ede des boutons directs pour $s$ et $r^2s$. Les sym\'etries d'axes diagonaux $rs$
et $r^3s$ s'obtiennent en composant une sym\'etrie directe avec un quart de tour.

\section{Les huit isom\'etries de Pentapol}

Le groupe di\'edral du carr\'e, not\'e ici $D_4$, est
\[
  \boxed{\;D_4=\{r^k,\ r^ks\mid k\in\Z/4\Z\}\;}.
\]
Il contient huit transformations :
\[
  D_4=\{\id,r,r^2,r^3,s,rs,r^2s,r^3s\}.
\]

\begin{center}
\begin{tabular}{clc}
\toprule
Transformation & Interpr\'etation & Minimum de boutons\\
\midrule
$\id$ & aucune transformation & $0$\\
$r$ & quart de tour trigonom\'etrique & $1$\\
$r^3=r^{-1}$ & quart de tour horaire & $1$\\
$s$ & sym\'etrie horizontale & $1$\\
$r^2s$ & sym\'etrie verticale & $1$\\
$r^2$ & demi-tour & $2$\\
$rs$ & sym\'etrie d'axe $y=x$ & $2$\\
$r^3s$ & sym\'etrie d'axe $y=-x$ & $2$\\
\bottomrule
\end{tabular}
\end{center}

Avec les quatre boutons de Pentapol, toute orientation est donc atteignable en au plus deux
appuis. Pour une pi\`ece qui poss\`ede ses propres sym\'etries, ce minimum peut encore diminuer.

\section{Calculer les compositions}

\subsection*{Deux rotations}

Les exposants s'additionnent modulo $4$ :
\[
  r^ar^b=r^{a+b\pmod 4}.
\]
Par exemple,
\[
  rr=r^2,qquad rr^3=r^4=\id,qquad r^3r^3=r^6=r^2.
\]

\subsection*{Rotation et sym\'etrie}

La relation fondamentale est
\[
  sr=r^{-1}s=r^3s.
\]
Elle implique
\[
  \boxed{\;srs=r^{-1}\;}.
\]
Une sym\'etrie inverse donc le sens d'une rotation. Visuellement, regarder une rotation dans un
miroir transforme le sens trigonom\'etrique en sens horaire.

\subsection*{Deux sym\'etries}

Dans Pentapol, la sym\'etrie horizontale est $s$ et la sym\'etrie verticale est $r^2s$. Une
sym\'etrie horizontale suivie d'une sym\'etrie verticale vaut donc
\[
  (r^2s)\circ s=r^2s^2=r^2.
\]
Le r\'esultat est une rotation de $\pi$, et non de $\pi/2$.

\section{Exemple num\'erique}

Prenons le point $A(2,1)$, d'affixe $z=2+i$.

\begin{center}
\begin{tabular}{ccc}
\toprule
Transformation & Calcul & Image de $A$\\
\midrule
$r$ & $i(2+i)=-1+2i$ & $(-1,2)$\\
$r^2$ & $-(2+i)$ & $(-2,-1)$\\
$r^3$ & $-i(2+i)=1-2i$ & $(1,-2)$\\
$s$ & $\conj{(2+i)}=2-i$ & $(2,-1)$\\
$rs$ & $i(2-i)=1+2i$ & $(1,2)$\\
$r^2s$ & $-(2-i)=-2+i$ & $(-2,1)$\\
$r^3s$ & $-i(2-i)=-1-2i$ & $(-1,-2)$\\
\bottomrule
\end{tabular}
\end{center}

Pour transformer une pi\`ece, on applique la m\^eme formule aux cinq cellules, toutes rep\'er\'ees
par rapport au m\^eme centre.

\section{Pourquoi aucune rotation de $\pi/4$ n'appara\^it}

Une rotation de $\pi/4$ aurait pour facteur
\[
  e^{i\pi/4}=\cos\frac{\pi}{4}+i\sin\frac{\pi}{4}
  =\frac{1+i}{\sqrt2}.
\]
Ce nombre n'appartient pas au cycle $\{1,i,-1,-i\}$. Il introduit g\'en\'eralement des coordonn\'ees
divis\'ees par $\sqrt2$, qui ne retombent pas sur la grille carr\'ee. Toute composition des
transformations de $D_4$ reste dans $D_4$ : aucune suite de boutons Pentapol ne peut donc produire
une rotation de $\pi/4$.

\section{Les pi\`eces n'ont pas toutes huit orientations visibles}

Le groupe $D_4$ contient huit transformations, mais les sym\'etries propres d'une pi\`ece peuvent
faire co\"incider plusieurs r\'esultats :

\begin{center}
\begin{tabular}{cc}
\toprule
Pi\`eces & Nombre d'orientations visibles\\
\midrule
X & $1$\\
I & $2$\\
T, U, V, W, Z & $4$\\
P, F, Y, L, N & $8$\\
\bottomrule
\end{tabular}
\end{center}

Pentapol calcule donc la distance minimale pour la pi\`ece concern\'ee, et non seulement dans le
groupe abstrait.

\section{Exercices}

\begin{enumerate}
  \item Simplifier $r^{11}$.
  \item Identifier la rotation $e^{i5\pi/2}$.
  \item Simplifier $r^2r^3$.
  \item Calculer $s^2$.
  \item Calculer une sym\'etrie horizontale suivie d'une sym\'etrie verticale.
  \item Donner la formule complexe de la sym\'etrie d'axe $y=x$.
  \item Expliquer pourquoi $rs\neq sr$.
  \item Une suite de boutons Pentapol peut-elle produire $e^{i\pi/4}$ ?
\end{enumerate}

\subsection*{Corrig\'e}

\begin{enumerate}
  \item $r^{11}=r^3$, car $11\equiv3\pmod4$.
  \item $e^{i5\pi/2}=e^{i\pi/2}=i$ : un quart de tour trigonom\'etrique.
  \item $r^2r^3=r^5=r$.
  \item $s^2=\id$.
  \item $(r^2s)s=r^2$ : une rotation de $\pi$.
  \item $z'=i\conj z$, soit $(x,y)\mapsto(y,x)$.
  \item L'op\'eration de droite est appliqu\'ee en premier ; changer l'ordre change l'axe diagonal.
  \item Non, car toute composition reste de la forme $r^k$ ou $r^ks$, avec $k$ modulo $4$.
\end{enumerate}

\section*{Formulaire final}

\[
\boxed{
\begin{aligned}
r^k(z)&=e^{ik\pi/2}z=i^kz,\\
S_{\varphi}(z)&=e^{2i\varphi}\conj z,\\
D_4&=\{r^k,r^ks\mid k\in\Z/4\Z\},\\
r^4&=\id,\qquad s^2=\id,\qquad srs=r^{-1}.
\end{aligned}}
\]

\end{document}
